\documentclass[12pt]{article}
\usepackage[utf8]{inputenc}
\usepackage[T1]{fontenc}
\usepackage{hyperref}
\title{Northern Straits Salish Language Resources for K-12 Curriculum Development}
\author{Compiled by Claude Sonnet 5.5 (claude-sonnet-5-5) for Andr\'as Kornai}
\date{}
\begin{document}
\maketitle

\section{Introduction}
Northern Straits Salish (ISO 639-3 str; Glottocode stra1244; Salishan) is spoken in British Columbia and Washington State. The searches found material only for the Saanich dialect (SEN\'C\!O\!\'T\!EN). That is the best-documented variety, and it is not the whole language. This report rests on one search. The survey reference lists the language as dormant, but sources found here describe active documentation work on SEN\'C\!O\!\'T\!EN. Nothing was found for the other varieties, such as Lummi or Samish.

\section{Primary Educational Resources}
No school curriculum, primer or textbook was found in the search results.

\section{Linguistic and Reference Materials}
\begin{itemize}
\item Timothy Montler, \emph{SEN\'C\!O\!\'T\!EN: a Dictionary of the Saanich Language} (University of Washington Press, 2018). It is a two-way English/SEN\'C\!O\!\'T\!EN dictionary of about 1,520 pages with more than 12,000 entries, most with example sentences. It includes an affix appendix and a root index. A digital open-access edition is reported to be available from Montler's website. Catalogue record: \url{https://knowledgebase.arts.ubc.ca/?p=3210}
\item UBC RelLex record for the dictionary. \url{https://wiki.ubc.ca/Documentation:RelLex/SEN%C4%86O%C5%A6EN:_a_Dictionary_of_the_Saanich_Language}
\item CBC news report on the dictionary. \url{https://amp.cbc.ca/lite/story/1.4797281}
\item Review in the BC Review. \url{https://thebcreview.ca/?p=35233}
\item Review at BC BookLook. \url{https://bcbooklook.com/547-recording-the-saanich-language/}
\item Review at ABC BookWorld. \url{https://abcbookworld.com/?p=21490}
\item A grammar by Montler is reported to be in progress. It was not verified.
\end{itemize}

\section{Religious and Community Resources}
None identified in the searches.

\section{Digital and Audio Resources}
A review reports that Montler's Saanich website has a word list with audio links and a spelling converter. The site address was not seen in the search results, so none is given here.

\section{Curriculum Development Context}
The 2018 dictionary is a strong base for materials development. A curriculum would need community partners and speaker input, which this search did not identify. Teaching programmes were not found.

\end{document}